\begin{titlepage}
    \thispagestyle{empty}

    \definecolor{coverpaper}{RGB}{248, 246, 240}
    \definecolor{covernavy}{RGB}{26, 48, 88}
    \definecolor{covergold}{RGB}{186, 146, 76}
    \definecolor{coverline}{RGB}{196, 202, 212}
    \definecolor{coverbrick}{RGB}{174, 96, 74}
    \definecolor{coverteal}{RGB}{82, 132, 136}
    \definecolor{coverplum}{RGB}{104, 97, 142}

    \begin{tikzpicture}[remember picture, overlay]
        \fill[coverpaper] (current page.south west) rectangle (current page.north east);
        \draw[covergold, line width=1.1pt]
            ([yshift=-1.0cm]current page.north west) --
            ([yshift=-1.0cm]current page.north east);
        \fill[covernavy] (current page.south west) rectangle ([yshift=2.65cm]current page.south east);

        % 统一主视觉：复平面与极坐标向量
        \begin{scope}[shift={([xshift=-4.9cm, yshift=7.6cm]current page.south east)}]
            \draw[coverline!82, line width=0.55pt] (-4.8, 0) -- (4.7, 0);
            \draw[coverline!82, line width=0.55pt] (0, -4.4) -- (0, 4.3);
            \draw[coverline!60, line width=0.45pt] (0, 0) circle (1.45);
            \draw[covergold!90, line width=1.05pt] (0, 0) circle (3.1);
            \draw[coverplum!44, line width=0.6pt, dash pattern=on 2.6pt off 2.0pt]
                (0, 0) circle (4.05);
            \fill[covergold!15] (0, 0) -- (1.08, 0)
                arc[start angle=0, end angle=58, radius=1.08] -- cycle;
            \draw[covergold!88!black, line width=0.85pt]
                (1.08, 0) arc[start angle=0, end angle=58, radius=1.08];
            \draw[
                covernavy,
                line width=1.18pt,
                -{Stealth[length=2.7mm, width=1.8mm]}
            ] (0, 0) -- (58:3.1);
            \node[text=covernavy, font=\fontsize{10.8pt}{12pt}\selectfont] at (58:3.7)
                {$re^{\mathrm{i}\theta}$};
            \node[text=covergold!88!black, font=\small] at (0.95, 0.6) {$\theta$};
            \node[text=black!42, font=\fontsize{9pt}{10pt}\selectfont] at (4.35, 0.28)
                {$\operatorname{Re}$};
            \node[text=black!42, font=\fontsize{9pt}{10pt}\selectfont] at (0.36, 4.0)
                {$\operatorname{Im}$};
        \end{scope}

        \node[
            anchor=west,
            text=black!42,
            font=\fontsize{10.2pt}{12pt}\selectfont
        ] at ([xshift=1.65cm, yshift=4.15cm]current page.south west)
            {$\displaystyle z=re^{\mathrm{i}\theta},\qquad \nabla^2 u + k^2 u = 0$};
    \end{tikzpicture}

    \vspace*{1.75cm}

    \begin{center}
        {\fontsize{18pt}{22pt}\selectfont\bfseries\color{covernavy!90}
            湖南师范大学\par}
        \vspace{0.12cm}
        {\fontsize{10.6pt}{13pt}\selectfont\color{black!42}
            Hunan Normal University\par}

        \vspace{1.55cm}

        {\fontsize{43pt}{49pt}\selectfont\bfseries\color{covernavy}
            数学物理方法\par}
        \vspace{0.35cm}
        {\fontsize{12.6pt}{15pt}\selectfont\itshape\color{covergold!88!black}
            Mathematical Methods for Physics\par}

        \vspace{0.5cm}

        {\color{covergold}
            \rule{0.17\textwidth}{0.8pt}\hspace{0.9em}
            {\large\color{covernavy!75}\ding{72}}\hspace{0.9em}
            \rule{0.17\textwidth}{0.8pt}
        \par}

        \vspace{0.62cm}

        {\fontsize{17.5pt}{22pt}\selectfont\color{covernavy!80}
            从物理问题出发的课程讲义\par}
        \vspace{0.28cm}
        {\fontsize{11pt}{14pt}\selectfont\color{black!48}
            复变函数、积分变换、数理方程与特殊函数\par}

        \vspace{0.9cm}

        \begin{tikzpicture}
            \foreach \x/\col/\txt in {
                -4.35/coverbrick/复变函数,
                -1.45/coverteal/积分变换,
                1.65/covergold/数理方程,
                4.55/coverplum/特殊函数
            }{
                \node[
                    rounded corners=1.8pt,
                    draw=\col!68,
                    fill=\col!10,
                    minimum width=2.45cm,
                    minimum height=0.7cm,
                    text=covernavy!88,
                    font=\fontsize{10pt}{12pt}\selectfont
                ] at (\x, 0) {\txt};
            }
        \end{tikzpicture}

        \vspace{0.4cm}

        {\fontsize{10.5pt}{12.8pt}\selectfont\color{black!48}
            面向物理学、光电信息科学等专业本科生的自学与教学\par}

        \vspace{1.05cm}

        {\fontsize{17pt}{20pt}\selectfont\bfseries\color{covernavy}
            陈祖成\par}
        \vspace{0.08cm}
        {\fontsize{10.4pt}{12pt}\selectfont\color{black!48}
            编著\par}
    \end{center}

    \begin{tikzpicture}[remember picture, overlay]
        \node[
            anchor=center,
            text=white,
            font=\fontsize{13pt}{16pt}\selectfont\bfseries
        ] at ([yshift=1.95cm]current page.south) {物理与电子科学学院};
        \node[
            anchor=center,
            text=white!75,
            font=\fontsize{9.7pt}{11.2pt}\selectfont
        ] at ([yshift=1.35cm]current page.south) {School of Physics and Electronics};
        \node[
            anchor=center,
            text=covergold!96,
            font=\fontsize{10.4pt}{12pt}\selectfont
        ] at ([yshift=0.68cm]current page.south) {2026 春};
    \end{tikzpicture}
\end{titlepage}
